\documentclass{article}
\usepackage[utf8]{inputenc}
\usepackage{geometry}
\usepackage{pdfpages}
\usepackage{mathtools}
\usepackage{array}
\usepackage{filecontents}
\usepackage{nicematrix}
\usepackage{bodeplot}
\graphicspath{{images/}}
\geometry{
	a4paper,
	total={170mm,257mm},
	left=20mm,
	top=20mm,
}
\usepackage{graphicx}
\usepackage{titling}
\usepackage[american,RPvoltages]{circuiTikz}
\usepackage{tikz, tikz-cd}
\usepackage{pgfplots}
\usepackage{siunitx}
\usepackage{float}
\usepackage{enumitem}
\usepackage{amsmath,amsfonts,stmaryrd,amssymb} % Math packages
\usepackage{schemabloc}
\usepackage{tcolorbox}
\usepackage{xcolor}
\usetikzlibrary{positioning}
\usepackage[scr]{rsfso}
\newcommand*{\Scale}[2][4]{\scalebox{#1}{$#2$}}%
\newcommand*{\Resize}[2]{\resizebox{#1}{!}{$#2$}}%

\newcommand{\Laplace}{\mathscr{L}}

\title{Análisis de circuito con OpAmp}
\author{Belmont}
\date{Agosto 2024}

\usepackage{fancyhdr}
\fancypagestyle{plain}{%  the preset of fancyhdr 
	\fancyhf{} % clear all header and footer fields
	\fancyfoot[L]{\thedate}
	\fancyhead[L]{Análisis de circuito con OpAmp}
	\fancyhead[R]{\theauthor}
}
\makeatletter
\def\@maketitle{%
	\newpage
	\null
	\vskip 1em%
	\begin{center}%
		\let \footnote \thanks
		{\LARGE \@title \par}%
		\vskip 1em%
		%{\large \@date}%
	\end{center}%
	\par
	\vskip 1em}
\makeatother

\usepackage{lipsum}  
\usepackage{cmbright}
\renewcommand{\figurename}{Figura}
%---------------------------------------------------
\tikzset{flipflop myRS/.style={flipflop,
		flipflop def={t1=R, t3=S, t6=Q, t4={\ctikztextnot{Q}}, tu={\ctikztextnot{RST}}, nu=1, n4=1 }}
}

\tikzset{flipflop myFFRS/.style={flipflop,
		flipflop def={t1=R, t3=S, t6=Q, t4={\ctikztextnot{Q}}, n4=1 }}
}
%---------------------------------------------------

\makeatletter
\newcommand*{\underarrow}{\def\@underarrow{\relax}\@ifstar{\@@underarrow}{\def\@underarrow{\hidewidth}\@@underarrow}}
\newcommand*{\@@underarrow}[2][]{\underset{\@underarrow\substack{\uparrow\if\relax\detokenize{#1}\relax\else\\#1\fi}\@underarrow}{#2}}

\newcommand*{\overarrow}{\def\@overarrow{\relax}\@ifstar{\@@overarrow}{\def\@overarrow{\hidewidth}\@@overarrow}}
\newcommand*{\@@overarrow}[2][]{\overset{\@overarrow\substack{\if\relax\detokenize{#1}\relax\else#1\\\fi\downarrow}\@overarrow}{#2}}
\makeatother

\makeatletter
\newcounter{elimination@steps}
\newcolumntype{R}[1]{>{\raggedleft\arraybackslash$}p{#1}<{$}}
\def\elimination@num@rights{}
\def\elimination@num@variables{}
\def\elimination@col@width{}
\newenvironment{elimination}[4][0]
{
	\setcounter{elimination@steps}{0}
	\def\elimination@num@rights{#1}
	\def\elimination@num@variables{#2}
	\def\elimination@col@width{#3}
	\renewcommand{\arraystretch}{#4}
	\start@align\@ne\st@rredtrue\m@ne
}
{
	\endalign
	\ignorespacesafterend
}
\newcommand{\eliminationstep}[2]
{
	\ifnum\value{elimination@steps}>0\sim\quad\fi
	\left[
	\ifnum\elimination@num@rights>0
	\begin{array}
		{@{}*{\elimination@num@variables}{R{\elimination@col@width}}
			|@{}*{\elimination@num@rights}{R{\elimination@col@width}}}
		\else
		\begin{array}
			{@{}*{\elimination@num@variables}{R{\elimination@col@width}}}
			\fi
			#1
		\end{array}
		\right]
		& 
		\begin{array}{l}
			#2
		\end{array}
		&%                                    moved second & here
		\addtocounter{elimination@steps}{1}
	}
	\makeatother
	
	\makeatletter
	\renewcommand*\env@matrix[1][*\c@MaxMatrixCols c]{%
		\hskip -\arraycolsep
		\let\@ifnextchar\new@ifnextchar
		\array{#1}}
	\makeatother
	
\setlength\parindent{0pt}
\begin{document}
	
	\maketitle
	
	\noindent\begin{tabular}{@{}ll}
		Author & \theauthor\\

	\end{tabular}

	
	\section*{Análisis}
Se el circuito con amplificador operacional de la figura 1. Determine el voltaje de salida $V_o$

\begin{figure}[H]
	\centering	
	\begin{tikzpicture}[scale=0.8]
		%\ctikzset{resistor = european}
		\ctikzset{bipoles/thickness=4}
		\ctikzset{monopoles/vcc/arrow={Stealth[width=6pt,length=9pt]}}
		\ctikzset{monopoles/vee/arrow={Stealth[width=6pt,length=9pt]}}
		\draw (0,0) node[op amp, anchor=+](OA1){$OA_1$};
		\draw(OA1.up) to[short] ++(0,0.4) node[vcc](VCC){$+V_{CC}$};
		\draw(OA1.down) to[short] ++(0,-0.4) node[vee](VEE){$-V_{EE}$};
		\draw (OA1.-) to[short] ++(-1,0) coordinate (p1) to[short] ++(0,2.5) coordinate (p2);
		\draw (p2) to[R, l=$R_3$] (p2 -| OA1.out) to[short] (OA1.out);
		\draw (p1) to[R, *-o, l_=$R_1$] ++(-3,0) node[left]{$V_{I_2}$};
		\draw (OA1.+) to[short] ++(-1,0) coordinate(p3) to[R, l=$R_2$, -o] ++(-3,0) node[left]{$V_{I_1}$};
		\draw (p3) to[R, *-*, l=$R_4$] ++(0,-3) coordinate(p4) to[R, l=$R_5$] ++(0,-3) node[ground]{};
		\draw (p4) to[R, l=$R_6$] ++(4,0) node[potentiometershape, anchor=wiper, rotate=90](P1){};
		\draw (P1.left) to[short] ++(0,-0.5) coordinate (vl);
		\draw (vl) node[vee](VEE){$V_{EE}$};
		\draw (P1.right) to[short] ++(0,0.5) coordinate (vu);
		\draw (vu) node[vcc](VCC){$V_{CC}$};
		\draw (P1.wiper) node[right=0.8cm]{$R_7$};
		
		\draw (OA1.out) node[below]{$V_{o_1}$} to[R, *-*, l=$R_8$] ++(4,0) coordinate (p6);
		\draw (p6) to[short] ++(0,-2.5) to[short] ++(1,0) node[op amp, anchor=-](OA2){};
		\draw (p6) to[C, l=$C$] (p6 -| OA2.out) to[short] (OA2.out);
		\draw (p6) to[R, l=$R_9$] ++(0,3) to[short] ++(3,0) node[potentiometershape, anchor=wiper, rotate=90](P2){};
		\draw (P2.left) to[short] ++(0,-0.5) coordinate (g1) node[ground]{};
		\draw (P2.right) to[short] ++(0,0.5) to[short] ++(2,0) coordinate (p7);
		\draw (p7) to[D, *-, l=$D$] (p7 |- g1) node[ground]{};
		\draw (p7) to[R, l=$R_{11}$] ++(3,0) node[vcc, rotate=-90](VCC){};
		\draw (P2.wiper) node[right=0.8cm]{$R_{10}$};
		\draw (P2.wiper) node[xshift=5.3cm, yshift=1cm]{$V_{CC}$};
		\draw (OA2.+) to[short] ++(-1,0) to[short] ++(0,-1) node[ground]{};
		\draw (OA2.out) to[short, *-o] ++(1,0) node[right]{$V_o$};
		\draw(OA2.up) to[short] ++(0,0.4) node[vcc](VCC){$+V_{CC}$};
		\draw(OA2.down) to[short] ++(0,-0.4) node[vee](VEE){$-V_{EE}$};
 		\end{tikzpicture}
		\caption{Circuito com amplificador operacional}	
		\label{fig:figura1}
	\end{figure}

Analizaremos el circuito por partes, nos enfocamos en la primera etapa que comprende al amplificador $OA_1$, notamos que en esencia se trata de un amplificador diferencial, utilizaremos el principio de superposición para encontrar el voltaje $V_{o_1}$.

Antes de proceder con el análisis modelaremos el potenciómetro con su circuito equivalente. Para modelar un potenciómetro, son necesarios dos parámetros: $R$ y $\alpha$. El parámetro $R$ especifica la resistencia total del potenciómetro ($R > 0$). El parámetro $\alpha$ representa la posición del contacto deslizante (la perilla del potenciómetro) y toma valores en el rango de $0 \leq \alpha \leq 1$. Los valores $\alpha = 0$ y $\alpha = 1$ corresponden a las posiciones extremas del contacto deslizante.
\begin{figure}[H]
	\centering	
	\begin{tikzpicture}[scale=0.8]
		%\ctikzset{resistor = european}
		\ctikzset{bipoles/thickness=4}
		\ctikzset{monopoles/vcc/arrow={Stealth[width=6pt,length=9pt]}}
		\ctikzset{monopoles/vee/arrow={Stealth[width=6pt,length=9pt]}}
		\draw (0,0) node[op amp, anchor=+](OA1){$OA_1$};
		\draw(OA1.up) to[short] ++(0,0.4) node[vcc](VCC){$+V_{CC}$};
		\draw(OA1.down) to[short] ++(0,-0.4) node[vee](VEE){$-V_{EE}$};
		\draw (OA1.-) to[short] ++(-1,0) coordinate (p1) to[short] ++(0,2.5) coordinate (p2);
		\draw (p2) to[R, l=$R_3$] (p2 -| OA1.out) to[short] (OA1.out);
		\draw (p1) to[R, *-o, l_=$R_1$] ++(-3,0) node[left]{$V_{I_2}$};
		\draw (OA1.+) to[short] ++(-1,0) coordinate(p3) to[R, l=$R_2$, -o] ++(-3,0) node[left]{$V_{I_1}$};
		\draw (p3) to[R, *-*, l=$R_4$] ++(0,-3) coordinate(p4) to[R, l=$R_5$] ++(0,-3) node[ground]{};
		\draw (p4) to[R, l=$R_6$] ++(4,0) node[potentiometershape, anchor=wiper, rotate=90](P1){};
		\draw (P1.left) to[short] ++(0,-0.5) coordinate (vl);
		\draw (vl) node[vee](VEE){$V_{EE}$};
		\draw (P1.right) to[short] ++(0,0.5) coordinate (vu);
		\draw (vu) node[vcc](VCC){$V_{CC}$};
		\draw (P1.wiper) node[right=0.8cm]{$R_7$};
	\end{tikzpicture}
	\caption{}	
	\label{fig:figura2}
\end{figure}

La figura 3 muestra el modelo para el potenciómetro que consiste de dos resistores.
\begin{figure}[H]
	\centering	
	\begin{tikzpicture}[scale=0.8]
		\ctikzset{bipoles/thickness=4}
		\draw (0,0) to[R, o-*, l_=$\alpha R_p$] ++(0,-3) coordinate (p1)
		to[R, l_=$(1-\alpha) R_p$, -o] ++(0,-3);
		\draw (p1) to[short, -o] ++(-1,0);
	\end{tikzpicture}
	\caption{Circuito equivalente de un potenciómetro}	
	\label{fig:figura3}
\end{figure}


Reemplazamos este modelo equivalente del potenciómetro en el circuito de la figura 2. El nuevo circuito se muestra en al figura 4. Ahora podemos proceder a aplicar el equivalente de Thévenin un par de veces para poder simplificar el circuito aún más.

Primero, encontramos el circuito equivalente de Thévenin sobre el nodo denomiando $V_{th_1}$ (el circuito se muestra en la figura 5). Para encontrar $R_{th_1}$ cortocircuitamos las fuentes de tensión para obtener:

\begin{align}
	R_{th_1} = \alpha R_7 \parallel (1 - \alpha) R_7
\end{align}

Al resolver obtenemos
\begin{align}
	R_{th_1} &= \frac{\alpha R_7 (1 - \alpha) R_7}{\alpha R_7 + (1 - \alpha) R_7}\nonumber \\
	&= \frac{\alpha R_7 (R_7 - \alpha R_7)}{\alpha R_7 + R_7 - \alpha R_7} \nonumber
\end{align}

\begin{tcolorbox}[colback=red!5!white,colframe=red!75!black,arc=0mm,boxrule=0mm,width=(\linewidth+0.5cm)]
	\begin{equation}
		R_{th_1} = \alpha R_7 (1 - \alpha)
	\end{equation}
\end{tcolorbox}

\begin{figure}[H]
	\centering	
	\begin{tikzpicture}[scale=0.8]
		%\ctikzset{resistor = european}
		\ctikzset{bipoles/thickness=4}
		\ctikzset{monopoles/vcc/arrow={Stealth[width=6pt,length=9pt]}}
		\ctikzset{monopoles/vee/arrow={Stealth[width=6pt,length=9pt]}}
		\draw (0,0) node[op amp, anchor=+](OA1){$OA_1$};
		\draw(OA1.up) to[short] ++(0,0.4) node[vcc](VCC){$+V_{CC}$};
		\draw(OA1.down) to[short] ++(0,-0.4) node[vee](VEE){$-V_{EE}$};
		\draw (OA1.-) to[short] ++(-1,0) coordinate (p1) to[short] ++(0,2.5) coordinate (p2);
		\draw (p2) to[R, l=$R_3$] (p2 -| OA1.out) to[short] (OA1.out);
		\draw (p1) to[R, *-o, l_=$R_1$] ++(-3,0) node[left]{$V_{I_2}$};
		\draw (OA1.+) to[short] ++(-1,0) coordinate(p3) to[R, l=$R_2$, -o] ++(-3,0) node[left]{$V_{I_1}$};
		\draw (p3) to[R, *-*, l=$R_4$] ++(0,-3) coordinate(p4) to[R, l=$R_5$] ++(0,-3) node[ground]{};
		\draw (p4) node[above left]{$V_a$} to[R, l=$R_6$] ++(5,0) node[above left]{$V_{th_1}$} coordinate (g2) to[R, l_=$\alpha R_7$] ++(0,3) node[vcc](VCC){$V_{CC}$};
		\draw (g2) to[R,*-, l=$(1-\alpha) R_7$] ++(0,-3) node[vee](VEE){$V_{EE}$};
	\end{tikzpicture}
	\caption{}	
	\label{fig:figura4}
\end{figure}

\begin{figure}[H]
	\centering	
	\begin{tikzpicture}[scale=0.8]
		\ctikzset{bipoles/thickness=4}
		\ctikzset{monopoles/vcc/arrow={Stealth[width=6pt,length=9pt]}}
		\ctikzset{monopoles/vee/arrow={Stealth[width=6pt,length=9pt]}}
		\draw (0,0) to[R, -*, l_=$\alpha R_7$, v^=$$] ++(0,-3) coordinate (p1)
		to[R, l_=$(1-\alpha) R_7$, v^=$$] ++(0,-3) coordinate (p3);
		\draw (p1) to[short, -o] ++(-3,0) coordinate (vth) node[left]{$$};
		\draw (0,0) to[short] ++(4,0) coordinate (p2) to[V, invert, l=$V_{CC}$] (p2 |- p1) coordinate (m);
		\draw (p3) to[short] ++(4,0) to[V, l_=$V_{EE}$] (m);
		\draw (m) to[short, *-] ++(1,0) to[short] ++(0,-0.2) node[ground]{};
		\draw[->]   (1.7,-2.3) arc(110:350:8mm) node[midway, xshift=5mm, yshift=5mm, font=\footnotesize] {$I_1$};
		\draw (vth) to[open, v=$V_{th_1}$] ++(0,-3) coordinate(vth1n);
		\draw (vth1n) to[short, o-] ++(0,0) node[ground]{};
	\end{tikzpicture}
	\caption{}	
	\label{fig:figura5}
\end{figure}

Empleamos la LTK sobre la malla del circuito de la figura 5:

\begin{align}
	V_{EE} + V_{CC} - \alpha R_7 I - (1- \alpha) R_7 I = 0 \nonumber \\
	V_{EE} + V_{CC} = [\alpha R_7 + (1- \alpha) R_7 ] I \nonumber \\
	V_{EE} + V_{CC} = [\alpha R_7 + R_7 - \alpha R_7 ] I \nonumber
\end{align}

\begin{align}
	I = \frac{V_{EE} + V_{CC}}{R_7}
\end{align}

Ahora encontramos el voltaje $V_{th_1}$ con respecto a tierra como se muestra a continuación
\begin{align}
	V_{th_1} - (1 - \alpha) R_7I + V_{EE} = 0
\end{align}

Reemplazamos (3) en (4) para obtener
\begin{align}
	V_{th_1} &= (1 - \alpha) \frac{R_7(V_{EE} + V_{CC})}{R7} - V_{EE} \nonumber \\
	&= (1 - \alpha) (V_{EE} + V_{CC}) - V_{EE} \nonumber \\
	&= V_{CC} + V_{EE} - \alpha V_{EE} - \alpha V_{CC} - V_{EE} \nonumber
\end{align}

\begin{tcolorbox}[colback=red!5!white,colframe=red!75!black,arc=0mm,boxrule=0mm,width=(\linewidth+0.5cm)]
	\begin{equation}
		V_{th_1} = V_{CC} - \alpha (V_{CC} + V_{EE})
	\end{equation}
\end{tcolorbox}

Redibujamos el circuito una vez obtenido el primer equivalente de Thévenin
\begin{figure}[H]
	\centering	
	\begin{tikzpicture}[scale=0.8]
		%\ctikzset{resistor = european}
		\ctikzset{bipoles/thickness=4}
		\ctikzset{monopoles/vcc/arrow={Stealth[width=6pt,length=9pt]}}
		\ctikzset{monopoles/vee/arrow={Stealth[width=6pt,length=9pt]}}
		\draw (0,0) node[op amp, anchor=+](OA1){$OA_1$};
		\draw(OA1.up) to[short] ++(0,0.4) node[vcc](VCC){$+V_{CC}$};
		\draw(OA1.down) to[short] ++(0,-0.4) node[vee](VEE){$-V_{EE}$};
		\draw (OA1.-) to[short] ++(-1,0) coordinate (p1) to[short] ++(0,2.5) coordinate (p2);
		\draw (p2) to[R, l=$R_3$] (p2 -| OA1.out) to[short] (OA1.out);
		\draw (p1) to[R, *-o, l_=$R_1$] ++(-3,0) node[left]{$V_{I_2}$};
		\draw (OA1.+) to[short] ++(-1,0) coordinate(p3) to[R, l=$R_2$, -o] ++(-3,0) node[left]{$V_{I_1}$};
		\draw (p3) to[R, *-*, l=$R_4$] ++(0,-3) coordinate(p4) to[R, l=$R_5$] ++(0,-3) node[ground]{};
		\draw (p4) node[above left]{$V_a$} to[R, l=$R_6$] ++(3,0) to[R, l=$R_{th_1}$] ++(3,0) to[V, invert, l=$V_{th_1}$] ++(0,-3) node[ground]{}; 
	\end{tikzpicture}
	\caption{}	
	\label{fig:figura6}
\end{figure}

Ahora encontramos el equivalente de Thévenin sobre el nodo $V_a$ como se muestra en la figura 7.

\begin{figure}[H]
	\centering	
	\begin{tikzpicture}[scale=0.8]
		%\ctikzset{resistor = european}
		\ctikzset{bipoles/thickness=4}
		\ctikzset{monopoles/vcc/arrow={Stealth[width=6pt,length=9pt]}}
		\ctikzset{monopoles/vee/arrow={Stealth[width=6pt,length=9pt]}}
		\draw (0,0) node[ground]{} to[V, l_=$V_{th_1}$] ++(0,3) to[R,l_=$R_{th_1}$] ++(-3,0) to[R, l_=$R_6$] coordinate (p1) ++(-3,0) to[R, l=$R_5$, *-] ++(0,-3) node[ground]{};
		\draw (p1) to[short, -o] ++(-2,0) node[left]{$V_a$};
		\draw[->]   (-3.3,2) arc(110:350:8mm) node[midway, xshift=5mm, yshift=5mm, font=\footnotesize] {$I_1$};
	\end{tikzpicture}
	\caption{}	
	\label{fig:figura7}
\end{figure}

De la figura 7 es fácil observar que la resistencia equivalente de Thévenin está dada por

\begin{tcolorbox}[colback=red!5!white,colframe=red!75!black,arc=0mm,boxrule=0mm,width=(\linewidth+0.5cm)]
	\begin{equation}
		R_{th_2} = R_5 \parallel (R_6 + R_{th_1}) = \frac{R5(R6 + R_{{th_1}})}{R_5 + R_6 + R_{th_1}}
	\end{equation}
\end{tcolorbox}

Mediante división de voltaje, obtenemos el voltaje $V_a$ como se muestra a continuación

\begin{tcolorbox}[colback=red!5!white,colframe=red!75!black,arc=0mm,boxrule=0mm,width=(\linewidth+0.5cm)]
\begin{equation}
	V_a = \frac{R_5}{R_5 + R_6 + R_{th_1}}V_{th_1} = V_{th_2}
	\end{equation}
\end{tcolorbox}

Esto nos posibilita redibujar el circuito de la figura 6 como se muestra a continuación:

\begin{figure}[H]
	\centering	
	\begin{tikzpicture}[scale=0.8]
		%\ctikzset{resistor = european}
		\ctikzset{bipoles/thickness=4}
		\ctikzset{monopoles/vcc/arrow={Stealth[width=6pt,length=9pt]}}
		\ctikzset{monopoles/vee/arrow={Stealth[width=6pt,length=9pt]}}
		\draw (0,0) node[op amp, anchor=+](OA1){$OA_1$};
		\draw(OA1.up) to[short] ++(0,0.4) node[vcc](VCC){$+V_{CC}$};
		\draw(OA1.down) to[short] ++(0,-0.4) node[vee](VEE){$-V_{EE}$};
		\draw (OA1.-) to[short] ++(-1,0) coordinate (p1) to[short] ++(0,2.5) coordinate (p2);
		\draw (p2) to[R, l=$R_3$] (p2 -| OA1.out) to[short] (OA1.out);
		\draw (p1) to[R, *-o, l_=$R_1$] ++(-3,0) node[left]{$V_{I_2}$};
		\draw (OA1.+) to[short] ++(-1,0) coordinate(p3) to[R, l=$R_2$, -o] ++(-3,0) node[left]{$V_{I_1}$};
		\draw (p3) to[R, *-, l=$R_4$] ++(0,-2.5) to[R, l=$R_{th_2}$] ++(0,-2.5) to[V, invert, l=$V_{th_2}$] ++(0,-2.5) node[ground]{};
		\draw (OA1.out) to[short, -o] ++(1,0) node[right]{$V_{o_1}$};
	\end{tikzpicture}
	\caption{}	
	\label{fig:figura8}
\end{figure}

Ya tenemos simplificada lo más posible la primera sección  del circuito, con esto en mente procedemos a obtener el voltaje de salida mediante superposición. 

Primero encontremos el voltaje de salida con respecto a $V_{I_2}$, esto se logra estableciendoa  cero la fuente de tensión $V_{I_1}$ y $V_{th_2}$

\begin{figure}[H]
	\centering	
	\begin{tikzpicture}[scale=0.8]
		%\ctikzset{resistor = european}
		\ctikzset{bipoles/thickness=4}
		\ctikzset{monopoles/vcc/arrow={Stealth[width=6pt,length=9pt]}}
		\ctikzset{monopoles/vee/arrow={Stealth[width=6pt,length=9pt]}}
		\draw (0,0) node[op amp, anchor=+](OA1){$OA_1$};
		\draw(OA1.up) to[short] ++(0,0.4) node[vcc](VCC){$+V_{CC}$};
		\draw(OA1.down) to[short] ++(0,-0.4) node[vee](VEE){$-V_{EE}$};
		\draw (OA1.-) to[short] ++(-1,0) coordinate (p1) to[short] ++(0,2.5) coordinate (p2);
		\draw (p2) to[R, l=$R_3$] (p2 -| OA1.out) to[short] (OA1.out);
		\draw (p1) to[R, *-o, l_=$R_1$] ++(-3,0) node[left]{$V_{I_2}$};
		\draw (OA1.+) to[short] ++(-1,0) coordinate(p3) to[R, l=$R_2$] ++(-3,0) node[ground]{};
		\draw (p3) to[R, *-, l=$R_4$] ++(0,-2.5) to[R, l=$R_{th_2}$] ++(0,-2.5) node[ground]{};
		\draw (OA1.out) to[short, -o] ++(1,0) node[right]{$V_{o_1}$};
	\end{tikzpicture}
	\caption{}	
	\label{fig:figura9}
\end{figure}

Del circuito de la figura 9, las resistencias $R_2$, $R_4$ y $R_{th_2}$ no tienen influencia sobre el voltaje de salida, por lo que fácilmente observamos que se trata de un amplificador inversor, así, el voltaje de salida está dado por

\begin{align}
	V_{o_1a} \bigg\rvert_{V_{I_1}, V_{th_2} = 0} = - \frac{R_3}{R_1} V_{I_2}
\end{align}

Ahora obtenemos el voltaje de salida con respecto a $V_{I_1}$, para esto desconectamos las fuentes $V_{I_2}$ y $V_{th_2}$ como se muestra en la figura 10.

\begin{figure}[H]
	\centering	
	\begin{tikzpicture}[scale=0.8]
		%\ctikzset{resistor = european}
		\ctikzset{bipoles/thickness=4}
		\ctikzset{monopoles/vcc/arrow={Stealth[width=6pt,length=9pt]}}
		\ctikzset{monopoles/vee/arrow={Stealth[width=6pt,length=9pt]}}
		\draw (0,0) node[op amp, anchor=+](OA1){$OA_1$};
		\draw(OA1.up) to[short] ++(0,0.4) node[vcc](VCC){$+V_{CC}$};
		\draw(OA1.down) to[short] ++(0,-0.4) node[vee](VEE){$-V_{EE}$};
		\draw (OA1.-) to[short] ++(-1,0) coordinate (p1) to[short] ++(0,2.5) coordinate (p2);
		\draw (p2) to[R, l=$R_3$] (p2 -| OA1.out) to[short] (OA1.out);
		\draw (p1) to[R, *-, l_=$R_1$] ++(-3,0) node[ground]{};
		\draw (OA1.+) to[short] ++(-1,0) coordinate(p3) to[R, l=$R_2$, -o] ++(-3,0) node[left]{$V_{I_1}$};
		\draw (p3) to[R, *-, l=$R_4$] ++(0,-2.5) to[R, l=$R_{th_2}$] ++(0,-2.5) node[ground]{};
		\draw (OA1.out) to[short, -o] ++(1,0) node[right]{$V_{o_1}$};
		\draw (p3) node[above]{$V_x$};
	\end{tikzpicture}
	\caption{}	
	\label{fig:figura10}
\end{figure}

Podemos notar que el circuito de la figura 10 se trata de un amplificador no inversor, antes de obtener el voltaje de salida, debemos encontrar el voltaje en el nodo $V_x$ el cual está dado por

\begin{figure}[H]
	\centering	
	\begin{tikzpicture}[scale=0.8]
		%\ctikzset{resistor = european}
		\ctikzset{bipoles/thickness=4}
		\ctikzset{monopoles/vcc/arrow={Stealth[width=6pt,length=9pt]}}
		\ctikzset{monopoles/vee/arrow={Stealth[width=6pt,length=9pt]}}
		\draw (0,0) to[V, l=$V_{I_1}$] ++(0,6) to[R, l=$R_2$] ++(4,0) coordinate (p1) to[R, l=$R_4$] ++(0,-3) to[R, l=$R_{th_2}$] ++(0,-3) to[short] ++(-4,0);
		\draw (p1) to[short, *-o] ++(2,0) node[right]{$V_x$};
	\end{tikzpicture}
	\caption{}	
	\label{fig:figura11}
\end{figure}

\begin{align}
	V_x = \frac{R_4 + R_{th_2}}{R_2 + R_4 + R_{th_2}}V_{I_1}
\end{align}

Podemos identificar fácilmente que el voltaje de salida está dada por

\begin{align}
	V_{o_1b} \bigg\rvert_{V_{I_2}, V_{th_2} = 0} =  \left( 1 + \frac{R_3}{R_1} \right) V_x = \left( 1 + \frac{R_3}{R_1} \right) \left( \frac{R_4 + R_{th_2}}{R_2 + R_4 + R_{th_2}} \right)V_{I_1}
\end{align}

Por último, calculamos el voltaje de salida con respecto a $V_{th_2}$, nuevamente, buscamos el voltaje $V_x$ para poder onbtener el voltaje de salida, sea el circuito de la figura 12, mediante división de tensión el voltaje $V_x$ es

\begin{figure}[H]
	\centering	
	\begin{tikzpicture}[scale=0.8]
		%\ctikzset{resistor = european}
		\ctikzset{bipoles/thickness=4}
		\ctikzset{monopoles/vcc/arrow={Stealth[width=6pt,length=9pt]}}
		\ctikzset{monopoles/vee/arrow={Stealth[width=6pt,length=9pt]}}
		\draw (0,0) node[op amp, anchor=+](OA1){$OA_1$};
		\draw(OA1.up) to[short] ++(0,0.4) node[vcc](VCC){$+V_{CC}$};
		\draw(OA1.down) to[short] ++(0,-0.4) node[vee](VEE){$-V_{EE}$};
		\draw (OA1.-) to[short] ++(-1,0) coordinate (p1) to[short] ++(0,2.5) coordinate (p2);
		\draw (p2) to[R, l=$R_3$] (p2 -| OA1.out) to[short] (OA1.out);
		\draw (p1) to[R, *-, l_=$R_1$] ++(-3,0) node[ground]{};
		\draw (OA1.+) to[short] ++(-1,0) coordinate(p3) to[R, l=$R_2$] ++(-3,0) node[ground]{};
		\draw (p3) to[R, *-, l=$R_4$] ++(0,-2.5) to[R, l=$R_{th_2}$] ++(0,-2.5) to[V, invert, l=$V_{th_2}$] ++(0,-2.5) node[ground]{};
		\draw (OA1.out) to[short, -o] ++(1,0) node[right]{$V_{o_1}$};
		\draw (p3) node[above]{$V_x$};
	\end{tikzpicture}
	\caption{}	
	\label{fig:figura12}
\end{figure}

\begin{align}
	V_x = \frac{R_2}{R_2 + R_4 + R_{th_2}}V_{th_2}
\end{align}

Así
\begin{align}
	V_{o_1c} \bigg\rvert_{V_{I_1}, V_{I_2} = 0} =  \left( 1 + \frac{R_3}{R_1} \right) V_x = \left( 1 + \frac{R_3}{R_1} \right) \left( \frac{R_2}{R_2 + R_4 + R_{th_2}} \right)V_{th_2}
\end{align}


Al combinar las ecuaciones (8), (10) y (12) encontramos el voltaje de salida completo, es decir


\begin{tcolorbox}[colback=red!5!white,colframe=red!75!black,arc=0mm,boxrule=0mm,width=(\linewidth+0.5cm)]
	\begin{equation}
		V_{o_1} = \left( 1 + \frac{R_3}{R_1} \right) \left( \frac{R_4 + R_{th_2}}{R_2 + R_4 + R_{th_2}} \right)V_{I_1} -\frac{R_3}{R_1}V_{I_2} + \left( 1 + \frac{R_3}{R_1} \right) \left( \frac{R_2}{R_2 + R_4 + R_{th_2}} \right)V_{th_2}
	\end{equation}
\end{tcolorbox}

O bien
\begin{tcolorbox}[colback=red!5!white,colframe=red!75!black,arc=0mm,boxrule=0mm,width=(\linewidth+0.5cm)]
	\begin{equation}
		V_{o_1} = \left( 1 + \frac{R_3}{R_1} \right) \left( \frac{1}{R_2 + R_4 + R_{th_2}} \right) \left[ (R4 + R_{th_2})V_{I_1} + R_2V_{th_2} \right] - \frac{R_3}{R_1} V_{I_2}
	\end{equation}
\end{tcolorbox}


Continuamos ahora con la segunda etapa del circuito, antes de proceder con el análisis, reemplazamos el potenciómetro $R_{10}$ por su modelo equivalente. También debemos notar la presencia del diodo $D_1$, éste establecerá un voltaje de $V_D$ volts sobre la terminal superior del potenciómetro. Por lo tanto, podemos redibujar el circuito de la figura 13 como se muestra en la figura 14.

\begin{figure}[H]
	\centering	
	\begin{tikzpicture}[scale=0.8]
		%\ctikzset{resistor = european}
		\ctikzset{bipoles/thickness=4}
		\ctikzset{monopoles/vcc/arrow={Stealth[width=6pt,length=9pt]}}
		\ctikzset{monopoles/vee/arrow={Stealth[width=6pt,length=9pt]}}
		
		\draw (0,0) node[left]{$V_{o_1}$} to[R, *-*, l=$R_8$] ++(4,0) coordinate (p6);
		\draw (p6) to[short] ++(0,-2.5) to[short] ++(1,0) node[op amp, anchor=-](OA2){};
		\draw (p6) to[C, l=$C$] (p6 -| OA2.out) to[short] (OA2.out);
		\draw (p6) to[R, l=$R_9$] ++(0,3) to[short] ++(3,0) node[potentiometershape, anchor=wiper, rotate=90](P2){};
		\draw (P2.left) to[short] ++(0,-0.5) coordinate (g1) node[ground]{};
		\draw (P2.right) to[short] ++(0,0.5) to[short] ++(2,0) coordinate (p7);
		\draw (p7) to[D, *-, l=$D_1$] (p7 |- g1) node[ground]{};
		\draw (p7) to[R, l=$R_{11}$] ++(3,0) node[vcc, rotate=-90](VCC){};
		\draw (P2.wiper) node[right=0.8cm]{$R_{10}$};
		\draw (P2.wiper) node[xshift=5.3cm, yshift=1cm]{$V_{CC}$};
		\draw (OA2.+) to[short] ++(-1,0) to[short] ++(0,-1) node[ground]{};
		\draw (OA2.out) to[short, *-o] ++(1,0) node[right]{$V_o$};
		\draw(OA2.up) to[short] ++(0,0.4) node[vcc](VCC){$+V_{CC}$};
		\draw(OA2.down) to[short] ++(0,-0.4) node[vee](VEE){$-V_{EE}$};
	\end{tikzpicture}
	\caption{}	
	\label{fig:figura13}
\end{figure}

\begin{figure}[H]
	\centering	
	\begin{tikzpicture}[scale=0.8]
		%\ctikzset{resistor = european}
		\ctikzset{bipoles/thickness=4}
		\ctikzset{monopoles/vcc/arrow={Stealth[width=6pt,length=9pt]}}
		\ctikzset{monopoles/vee/arrow={Stealth[width=6pt,length=9pt]}}
		
		\draw (0,0) node[left]{$V_{o_1}$} to[R, *-*, l=$R_8$, f=$I_{R8}$] ++(4,0) coordinate (p6);
		\draw (p6) to[short] ++(0,-2.5) to[short] ++(1,0) node[op amp, anchor=-](OA2){};
		\draw (p6) to[C, l=$C_1$, f>^=$I_{C1}$] (p6 -| OA2.out) to[short] (OA2.out);
		\draw (p6) to[R, l=$R_9$, f=$I_{R9}$] ++(0,3) to[short] ++(5,0) coordinate (h1);
		\draw (h1) node[above left]{$V_b$} to[R, l=$(1-\alpha_1)R_{10}$] ++(0,-3) node[ground]{};
		\draw (h1) to[R, *-, l_=$\alpha_1 R_{10}$] ++(0,3) to[short] ++(3,0) coordinate (h2) to[D, l=$D_1$] ++(0,-3) node[ground]{};
		\draw (h2) to[short, *-o] ++(2,0) node[right]{$V_D$};
		\draw (OA2.+) to[short] ++(-1,0) to[short] ++(0,-0.2) node[ground]{};
		\draw (OA2.out) to[short, *-o] ++(1,0) node[right]{$V_o$};
	\end{tikzpicture}
	\caption{Reemplazando el potenciómetro por su modelo equivalente.}	
	\label{fig:figura14}
\end{figure}

Calculamos el equivalente de Thévenin sobre el nodo $V_b$, modelamos además el diodo $D_1$ como una fuente de tensión con un voltaje de $V_D$ volts el cual alimentará al potenciómetro.
\begin{figure}[H]
	\centering	
	\begin{tikzpicture}[scale=0.8]
		%\ctikzset{resistor = european}
		\ctikzset{bipoles/thickness=4}
		\ctikzset{monopoles/vcc/arrow={Stealth[width=6pt,length=9pt]}}
		\ctikzset{monopoles/vee/arrow={Stealth[width=6pt,length=9pt]}}
		\draw (0,0) to[V, l=$V_{D}$] ++(0,6) to[short] ++(-4,0) to[R, l=$\alpha_1 R_{10}$] ++(0,-3) coordinate (p1) to[R, l=$(1-\alpha_1)R_{10}$] ++(0,-3) to[short] ++(4,0);
		\draw (p1) to[short, *-o] ++(-2,0) node[left]{$V_{th_3}$};
	\end{tikzpicture}
	\caption{}	
	\label{fig:figura15}
\end{figure}

Podemos calcular el voltaje de Thévenin mediante división de tensión:

\begin{tcolorbox}[colback=red!5!white,colframe=red!75!black,arc=0mm,boxrule=0mm,width=(\linewidth+0.5cm)]
	\begin{equation}
		V_{th_3} = \frac{(1 - \alpha_1) R_{10} V_D}{(1 - \alpha_1) R_{10} + \alpha_1 R_{10}} = (1 - \alpha_1)V_D
	\end{equation}
\end{tcolorbox}

La resistencia de Thévenin es similar a la de la ecuación (2), por lo cual 
\begin{tcolorbox}[colback=red!5!white,colframe=red!75!black,arc=0mm,boxrule=0mm,width=(\linewidth+0.5cm)]
	\begin{equation}
		R_{th_3} = \alpha_1 R_{10}(1 - \alpha_1)
	\end{equation}
\end{tcolorbox}

\begin{figure}[H]
	\centering	
	\begin{tikzpicture}[scale=0.8]
		%\ctikzset{resistor = european}
		\ctikzset{bipoles/thickness=4}
		\ctikzset{monopoles/vcc/arrow={Stealth[width=6pt,length=9pt]}}
		\ctikzset{monopoles/vee/arrow={Stealth[width=6pt,length=9pt]}}
		
		\draw (0,0) node[left]{$V_{o_1}$} to[R, *-*, l=$R_8$, f=$I_{R8}$] ++(4,0) coordinate (p6) node[below right]{$V_x$};
		\draw (p6) to[short] ++(0,-2.5) to[short] ++(1,0) node[op amp, anchor=-](OA2){};
		\draw (p6) to[C, l=$C_1$, f>^=$I_{C1}$] (p6 -| OA2.out) to[short] (OA2.out);
		\draw (p6) to[R, l=$R_9$, f<=$I_{R9}$] ++(0,3) to[R, -o, l=$R_{th_3}$] ++(4,0) coordinate (h1) node[right]{$V_{th_3}$};
		\draw (OA2.+) to[short] ++(-1,0) to[short] ++(0,-0.2) node[ground]{};
		\draw (OA2.out) to[short, *-o] ++(1,0) node[right]{$V_o$};
	\end{tikzpicture}
	\caption{}	
	\label{fig:figura16}
\end{figure}

Empleamos la LCK sobre el nodo $V_x$

\begin{align}
	I_{R8} + I_{R9} = I_{C1}
\end{align}

Mediante ley de Ohm:

\begin{align}
	\frac{V_{o_1} - V_x}{R_8} + \frac{V_{th_3} - V_x}{R_9 + R_{th_3}} = C\frac{d(V_x - V_o)}{dt}
\end{align}


Sin embargo, tenemos que $V_x = 0$, al reemplazarlo en la ecuación (18) y haciendo que $R_x = R_9 + R_{th_3}$ obtenemos

\begin{align}
	\frac{V_{o_1}}{R_8} + \frac{V_{th_3}}{R_x} = -C\frac{d(V_o)}{dt}
\end{align}

Despejamos $V_o$ de (19)

\begin{align}
	-\frac{1}{CR_8}V_{o1} - \frac{1}{CR_x}V_{th_3} = \frac{dV_o}{dt} \nonumber \\
	dV_o = - \left( \frac{1}{CR_8}V_{o1} + \frac{1}{CR_x}V_{th_3} \right)dt \nonumber \\
	\int dV_o = - \int \left( \frac{1}{CR_8}V_{o1} + \frac{1}{CR_x}V_{th_3} \right)dt
\end{align}

Es decir

\begin{tcolorbox}[colback=red!5!white,colframe=red!75!black,arc=0mm,boxrule=0mm,width=(\linewidth+0.5cm)]
	\begin{equation}
		V_o = - \frac{1}{C} \left( \frac{V_{o_1}}{R_8} + \frac{V_{th_3}}{R_x} \right) t
	\end{equation}
\end{tcolorbox}


\end{document}